\hypertarget{chapter-06-page-133}{}
\subsection{一个插值结果}
\label{sec:ch6-interpolation}

应用插值定理, 例如 Marcinkiewicz 插值定理时, 我们经常利用算子从 $L^\infty$ 到 $L^\infty$ 有界这一事实. 这里说明, 可以用更弱的从 $L^\infty$ 到 BMO 有界来代替它.

\begin{Theorem}
\label{thm:ch6-linfty-bmo-interpolation}
设线性算子 $T$ 在某个 $L^{p_0}$ 上有界, 其中 $1\le p_0<\infty$, 并且从 $L^\infty$ 到 BMO 有界. 则对每个 $p_0<p<\infty$, $T$ 在 $L^p$ 上有界.
\end{Theorem}

定理~\ref{thm:ch6-linfty-bmo-interpolation} 的证明需要两个引理.

\begin{Lemma}
\label{lem:ch6-dyadic-sharp-norm}
若 $1\le p_0\le p<\infty$ 且 $f\in L^{p_0}$, 则
\begin{equation}
\label{eq:ch6-dyadic-sharp-norm}
 \int_{\mathbb R^n}M_df(x)^p\,dx
 \le C\int_{\mathbb R^n}M^\#f(x)^p\,dx,
\end{equation}
其中 $M_d$ 是 \eqref{eq:ch2-dyadic-maximal} 定义的二进极大算子.
\end{Lemma}

由引理~\ref{lem:ch2-dyadic-to-cube-maximal}, 可以把 \eqref{eq:ch6-dyadic-sharp-norm} 中的二进极大函数换成 Hardy--Littlewood 极大函数. 因而, 尽管逐点不等式 $M^\#f(x)\le CMf(x)$ 的反向不等式一般不成立, 引理~\ref{lem:ch6-dyadic-sharp-norm} 给出了替代的范数不等式.

引理~\ref{lem:ch6-dyadic-sharp-norm} 的证明使用所谓的 good-$\lambda$ 不等式; 更确切地说, 它依赖以下结果.

\begin{Lemma}
\label{lem:ch6-good-lambda}
若 $f\in L^{p_0}$, 其中 $1\le p_0<\infty$, 则对所有 $\gamma>0$ 和 $\lambda>0$,
\[
 \bigl|\{x\in\mathbb R^n:M_df(x)>2\lambda,\ M^\#f(x)\le\gamma\lambda\}\bigr|
 \le2^n\gamma\bigl|\{x\in\mathbb R^n:M_df(x)>\lambda\}\bigr|.
\]
\end{Lemma}

\begin{Proof}
不失一般性, 可设 $f$ 非负. 固定 $\lambda,\gamma>0$. 按定理~\ref{thm:ch2-calderon-zygmund-decomposition}, 在高度 $\lambda$ 对 $f$ 作 Calderón--Zygmund 分解, 则集合 $\{x\in\mathbb R^n:M_df(x)>\lambda\}$ 可写成互不相交的极大二进立方体之并. 在相应证明中, 可用 $f\in L^{p_0}$ 代替 $f\in L^1$ 的假设. 若 $Q$ 是这些立方体之一, 只需证明
\begin{equation}
\label{eq:ch6-good-lambda-cube}
 |\{x\in Q:M_df(x)>2\lambda,\ M^\#f(x)\le\gamma\lambda\}|
 \le2^n\gamma|Q|.
\end{equation}

设 $\widetilde Q$ 是包含 $Q$ 且边长为其两倍的二进立方体. 由 $Q$ 的极大性, $f_{\widetilde Q}\le\lambda$. 此外, 若 $x\in Q$ 且 $M_df(x)>2\lambda$, 则 $M_d(f\chi_Q)(x)>2\lambda$, 所以对这些 $x$,
\[
 M_d((f-f_{\widetilde Q})\chi_Q)(x)
 \ge M_d(f\chi_Q)(x)-f_{\widetilde Q}>\lambda.
\]

\hypertarget{chapter-06-page-134}{}
由定理~\ref{thm:ch2-dyadic-maximal} 中 $M_d$ 的弱型 $(1,1)$ 不等式,
\[
\begin{aligned}
 |\{x\in Q:M_d((f-f_{\widetilde Q})\chi_Q)(x)>\lambda\}|
 &\le\frac1\lambda\int_Q|f(x)-f_{\widetilde Q}|\,dx\\
 &\le\frac{2^n|Q|}{\lambda}\frac1{|\widetilde Q|}
 \int_{\widetilde Q}|f(x)-f_{\widetilde Q}|\,dx\\
 &\le\frac{2^n|Q|}{\lambda}\inf_{x\in Q}M^\#f(x).
\end{aligned}
\]
若 \eqref{eq:ch6-good-lambda-cube} 左端的集合为空, 则无需证明. 否则对某个 $x\in Q$ 有 $M^\#f(x)\le\gamma\lambda$, 因而 \eqref{eq:ch6-good-lambda-cube} 成立.
\end{Proof}

\begin{Proof}[引理~\ref{lem:ch6-dyadic-sharp-norm} 的证明]
对 $N>0$, 令
\[
 I_N=\int_0^Np\lambda^{p-1}
 |\{x\in\mathbb R^n:M_df(x)>\lambda\}|\,d\lambda.
\]
$I_N$ 有限, 因为
\[
 I_N\le\frac p{p_0}N^{p-p_0}
 \int_0^Np_0\lambda^{p_0-1}
 |\{x\in\mathbb R^n:M_df(x)>\lambda\}|\,d\lambda,
\]
而 $f\in L^{p_0}$ 蕴含 $M_df\in L^{p_0}$. 进一步,
\[
\begin{aligned}
 I_N
 &=2^p\int_0^{N/2}p\lambda^{p-1}
 |\{x\in\mathbb R^n:M_df(x)>2\lambda\}|\,d\lambda\\
 &\le2^p\int_0^{N/2}p\lambda^{p-1}
 \bigl(|\{M_df>2\lambda,\ M^\#f\le\gamma\lambda\}|
 +|\{M^\#f>\gamma\lambda\}|\bigr)\,d\lambda\\
 &\le2^{p+n}\gamma I_N
 +\frac{2^p}{\gamma^p}\int_0^{\gamma N/2}p\lambda^{p-1}
 |\{x\in\mathbb R^n:M^\#f(x)>\lambda\}|\,d\lambda.
\end{aligned}
\]
固定 $\gamma$ 使 $2^{p+n}\gamma=1/2$, 则
\begin{equation}
\label{eq:ch6-truncated-distribution-bound}
 I_N\le\frac{2^{p+1}}{\gamma^p}\int_0^{\gamma N/2}p\lambda^{p-1}
 |\{x\in\mathbb R^n:M^\#f(x)>\lambda\}|\,d\lambda.
\end{equation}
若 $\int(M^\#f)^p=\infty$, 则无需证明. 若该积分有限, 在 \eqref{eq:ch6-truncated-distribution-bound} 中令 $N\to\infty$, 即得 \eqref{eq:ch6-dyadic-sharp-norm}.
\end{Proof}

\begin{Proof}[定理~\ref{thm:ch6-linfty-bmo-interpolation} 的证明]
复合算子 $M^\#\circ T$ 是次线性算子. 因为两个算子都在 $L^{p_0}$ 上有界, 所以该复合算子也在 $L^{p_0}$ 上有界; 它还在 $L^\infty$ 上有界, 因为
\[
 \|M^\#(Tf)\|_\infty=\|Tf\|_*\le C\|f\|_\infty.
\]
由 Marcinkiewicz 插值定理, $M^\#\circ T$ 在 $L^p$, $p_0<p<\infty$, 上有界.

\hypertarget{chapter-06-page-135}{}
现在设 $f\in L^p$ 且具有紧支集. 则 $f\in L^{p_0}$, 从而 $Tf\in L^{p_0}$. 因此可对 $Tf$ 应用引理~\ref{lem:ch6-dyadic-sharp-norm}. 特别地, 因为几乎处处有 $|Tf(x)|\le M_d(Tf)(x)$,
\[
\begin{aligned}
 \int_{\mathbb R^n}|Tf(x)|^p\,dx
 &\le\int_{\mathbb R^n}M_d(Tf)(x)^p\,dx\\
 &\le C\int_{\mathbb R^n}M^\#(Tf)(x)^p\,dx
 \le C\int_{\mathbb R^n}|f(x)|^p\,dx.
\end{aligned}
\]
\end{Proof}
